% large-rn-app-architecture.tex — structuring a large React Native codebase: workspaces +
% monorepo task runners (Nx / Turborepo), package layers and the dependency rule, feature
% packages with a public API ("exports"), enforcing boundaries with lint, code ownership,
% Metro in a monorepo (one react-native copy, package exports, autolinking), navigation
% assembled from features (React Navigation / Expo Router thin route files), state ownership
% per layer, universal apps (react-native-web + Expo Router web), a native-modules package,
% brownfield integration into existing native apps.
% Not repeated (neighbours): state taxonomy (react-native/rn-state-and-data.tex);
% navigator mechanics (react-native/rn-navigation.tex).
% Sources (checked 2026-09-25): docs.expo.dev/guides/monorepos (Metro auto-configured for
% monorepos since SDK 52; Bun/npm/pnpm/Yarn; pnpm isolated supported from SDK 54 else
% nodeLinker: hoisted; "Duplicate React Native versions in a single monorepo are not
% supported"; autolinkingModuleResolution SDK 54 opt-in, automatic SDK 55);
% reactnative.dev/blog/2025/04/08/react-native-0.79 (package.json "exports"/"imports"
% resolution on by default from RN 0.79 / Metro 0.82; opt out with
% resolver.unstable_enablePackageExports = false); reactnative.dev/docs/integration-with-existing-apps
% (RCTReactNativeFactory, ReactActivity/ReactHost, move the native app into ios/ + android/);
% nx.dev enforce-module-boundaries (@nx/enforce-module-boundaries, sourceTag,
% onlyDependOnLibsWithTags); docs.expo.dev/router/introduction (Android, iOS and web share one
% file-based navigation structure; static rendering on web); github.com/callstack/react-native-brownfield
% (helpers + CLI packaging RN as XCFramework / AAR).
% @labels: area=architecture kind=architecture platform=cross-platform topic=architecture,build,tooling discussion=84
% @tags: monorepo, workspaces, nx, turborepo, feature-modules, module-boundaries, codeowners, metro, package-exports, expo-router, react-native-web, brownfield
\documentclass[8pt]{extarticle}
\usepackage{printup-sheet}
\usepackage{array}

\lstdefinelanguage{TSSheet}{
  morekeywords={import,from,export,default,const,let,type,interface,function,return,
    as,true,false},
  sensitive=true, morecomment=[l]{//}, morestring=[b]", morestring=[b]',
  literate={=>}{{\hbox{=}\hbox{>}}}2 {->}{{\hbox{-}\hbox{>}}}2 {--}{{\hbox{-}\hbox{-}}}2
           {==}{{\hbox{=}\hbox{=}}}2 {!=}{{\hbox{!}\hbox{=}}}2 {**}{{\hbox{*}\hbox{*}}}2}
\lstset{basicstyle=\ttfamily\scriptsize, aboveskip=2pt, belowskip=2pt}

\newcommand\ct[1]{\texttt{#1}}
\newcommand\hd[1]{\par\vspace{3pt}\noindent{\bfseries\color{sheetBlue}#1}\par\vspace{1pt}}

\tikzset{
  pk/.style={box, font=\scriptsize, inner sep=1.5pt, minimum height=6mm, align=center},
  app/.style={pk, draw=sheetGrey, fill=black!5},
  ft/.style={pk, draw=sheetBlue, fill=sheetBlue!8},
  dt/.style={pk, draw=sheetGreen!60!black, fill=sheetGreen!10},
  ui/.style={pk, draw=sheetOrange, fill=sheetOrange!10},
  nt/.style={pk, draw=sheetBrown, fill=sheetBrown!10},
  lay/.style={font=\tiny\bfseries, text=sheetGrey, anchor=east},
  lbl/.style={font=\tiny, text=black!75, inner sep=1pt, align=center},
  ok/.style={->, thick, draw=sheetGrey},
  no/.style={->, thick, draw=sheetRed, dashed},
}

\begin{document}

\sheettitle{Large React Native app architecture}{architecture · folio}

\oneliner{A big RN codebase stays fast to change when it is a \textbf{monorepo of packages}
with \textbf{one dependency direction} — apps $\to$ features $\to$ data / UI $\to$ native
modules — each package exposing a \textbf{public API} (\ct{"exports"}), boundaries
\textbf{enforced by lint}, an \textbf{owner} per package, a \textbf{task graph} that builds and
tests only what changed, and exactly \textbf{one copy of \ct{react-native}}.}

\vspace{3pt}
\noindent\begin{tikzpicture}[sheet]
  % layers
  \node[lay] at (1.35,2.75) {apps};
  \node[lay] at (1.35,1.75) {features};
  \node[lay] at (1.35,0.75) {shared};
  \node[lay] at (1.35,-0.25) {platform};
  \node[app] (mob) at (2.8,2.75) {\ct{apps/mobile}\\\tiny Expo / RN shell};
  \node[app] (web) at (5.3,2.75) {\ct{apps/web}\\\tiny RN-web / Expo Router};
  \node[app] (sb) at (7.6,2.75) {\ct{apps/storybook}};
  \node[ft] (fa) at (2.8,1.75) {\ct{feature-feed}};
  \node[ft] (fb) at (5.3,1.75) {\ct{feature-checkout}};
  \node[ft] (fc) at (7.6,1.75) {\ct{feature-profile}};
  \node[dt] (api) at (2.8,0.75) {\ct{data}: API client +\\\tiny query keys, models};
  \node[ui] (ds) at (5.3,0.75) {\ct{ui}: design system};
  \node[dt] (core) at (7.6,0.75) {\ct{core}: auth, i18n,\\\tiny analytics facade};
  \node[nt] (nat) at (3.9,-0.25) {\ct{native-*}: Expo / Turbo modules};
  \node[nt, draw=sheetGrey, fill=white] (cfg) at (7.0,-0.25) {\ct{config}: eslint, tsconfig};
  \foreach \a/\b in {mob/fa,mob/fb,mob/fc,web/fa,web/fb,sb/ds,fa/api,fa/ds,fb/api,fb/ds,fb/core,fc/ds,fc/core,api/nat,core/nat}
    \draw[ok] (\a) -- (\b);
  \draw[no] (fb.east) -- node[lbl, above=1pt, text=sheetRed, font=\scriptsize\bfseries]{×} (fc.west);
  \draw[no] (ds.east) -- node[lbl, above=1pt, text=sheetRed, font=\scriptsize\bfseries]{×} (core.west);
  \node[lbl, anchor=west, text=sheetGrey] at (0.0,-0.85) {grey = allowed imports, always \textbf{downward}; \textcolor{sheetRed}{red dashed} (feature$\to$feature, ui$\to$core) fails CI};
  \draw[sheetGrey!50] (9.0,3.2) -- (9.0,-1.0);
  % affected graph
  \node[font=\bfseries\small, text=sheetBlue, anchor=west] at (9.15,3.05) {PR touches \ct{ui} $\to$ what runs?};
  \node[ui, very thick] (u2) at (10.4,2.25) {\ct{ui} changed};
  \node[ft, fill=sheetOrange!20] (a1) at (9.8,1.3) {feed};
  \node[ft, fill=sheetOrange!20] (a2) at (11.0,1.3) {checkout};
  \node[ft, fill=sheetOrange!20] (a3) at (12.2,1.3) {profile};
  \node[app, fill=sheetOrange!20] (a4) at (10.4,0.35) {mobile · web · storybook};
  \node[dt, fill=black!3, draw=sheetGrey] (c1) at (13.6,2.25) {\ct{data}};
  \node[nt, fill=black!3, draw=sheetGrey] (c2) at (13.6,1.3) {\ct{native-*}};
  \draw[hot] (u2) -- (a1); \draw[hot] (u2) -- (a2); \draw[hot] (u2) -- (a3);
  \draw[hot] (a1) -- (a4); \draw[hot] (a2) -- (a4);
  \node[lbl, anchor=west, align=left] at (13.0,0.55) {untouched: \textbf{cache hit}\\(content hash of inputs)};
  \node[lbl, anchor=west, align=left, text=sheetOrange] at (9.15,-0.35) {\ct{nx affected -t lint test typecheck} · Turborepo \ct{\hbox{-}\hbox{-}filter}\\= the \emph{reverse} dependency graph of the diff; remote cache shared by CI + laptops};
\end{tikzpicture}

\begin{multicols}{2}
\raggedright\footnotesize\setstretch{1.0}

\hd{The repo}
\begin{itemize}
  \item \textbf{Workspaces} (Yarn, pnpm, npm, Bun) symlink \ct{packages/*} into
        \ct{node\_modules}; \textbf{Nx} or \textbf{Turborepo} add a \textbf{task graph} from those
        dependencies, input hashing, local + \textbf{remote cache}, and ``affected'' runs.
  \item \textbf{One \ct{react-native}, one \ct{react}}: Expo — ``duplicate React Native versions
        in a single monorepo are not supported''. Pin them in the app, \ct{peerDependencies} in
        packages; check \ct{npm why react-native} / \ct{pnpm why}.
  \item \textbf{Metro}: Expo SDK 52+ configures a monorepo automatically; bare RN adds the root
        to \ct{watchFolders}. RN 0.79 (Metro 0.82) resolves \ct{package.json} \ct{"exports"}
        and \ct{"imports"} by default (opt out: \ct{resolver.unstable\_enablePackageExports}).
  \item \textbf{pnpm}: isolated \ct{node\_modules} catches \emph{phantom} (undeclared) deps;
        Expo supports it from SDK 54, else \ct{nodeLinker: hoisted}.
  \item \textbf{Native code autolinks from the app}: a native dependency must be resolvable from
        \ct{apps/mobile} (Expo SDK 54 \ct{autolinkingModuleResolution}, automatic in SDK 55).
\end{itemize}

\hd{Packages and boundaries}
\begin{itemize}
  \item \textbf{Feature package} = vertical slice: screens, its navigator, hooks, query keys,
        analytics events. Public surface = \ct{index.ts} + \ct{"exports"}; deep imports
        (\ct{feature-feed/src/x}) are not resolvable.
  \item \textbf{Features never import features.} Cross-feature needs go \emph{down} (a shared
        \ct{data}/\ct{core} package), \emph{through navigation} (route + params), or through an
        injected interface the app shell wires up.
  \item \textbf{Enforce}: Nx tags (\ct{type:feature}, \ct{type:ui}, \ct{scope:checkout}) +
        \ct{@nx/enforce-module-boundaries} \ct{depConstraints}; or \ct{no-restricted-imports},
        \ct{import/no-cycle}; TS project references for build order.
  \item \textbf{Ownership}: \ct{CODEOWNERS} per package path $\to$ the owning team must approve;
        a package without an owner is where rot starts.
  \item \textbf{State by owner}: server cache (TanStack Query) keys \emph{owned by} the data
        package — invalidation from any feature uses the same key factory; client stores
        (Zustand) \emph{per feature}; truly global = session, theme, flags only.
\end{itemize}

\hd{Example — public API, boundary rule, thin route}
\begin{lstlisting}[language=TSSheet]
// packages/feature-feed/package.json
{ "name": "@acme/feature-feed",
  "exports": { ".": "./src/index.ts" },           // nothing else
  "peerDependencies": { "react": "*", "react-native": "*" } }
// eslint.config.mjs (root)
'@nx/enforce-module-boundaries': ['error', { depConstraints: [
  { sourceTag: 'type:feature',
    onlyDependOnLibsWithTags: ['type:data', 'type:ui', 'type:core'] },
  { sourceTag: 'type:ui', onlyDependOnLibsWithTags: ['type:ui'] } ] }]
// apps/mobile/app/(tabs)/feed.tsx  -- Expo Router route file
export { FeedScreen as default } from '@acme/feature-feed';
\end{lstlisting}

\hd{Versioning, build, release}
\begin{itemize}
  \item Internal packages: \ct{"workspace:*"} — no versions, always head; only packages
        \emph{published} outside (an SDK, the design system for another repo) get semver +
        Changesets.
  \item \ct{tsc -b} with project references (or Nx \ct{typecheck}) per package — incremental,
        and a cycle is a build error.
  \item One native build per app; OTA channel per app + \ct{runtimeVersion}; a feature flag per
        risky feature so a bad feature is switched off, not hot-fixed.
  \item \textbf{Polyrepo} instead? Only when teams must release independently — the price is
        version skew, cross-repo PRs and a published package for every shared change.
  \item \textbf{Migrating a big app in}: start with \ct{config} + \ct{ui}, then carve one
        feature at a time out of \ct{apps/mobile/src} (strangler) — boundaries as
        \emph{warnings} first, errors once the count reaches zero.
\end{itemize}

\columnbreak

\hd{Navigation architecture}
\begin{itemize}
  \item The \textbf{app shell owns the root} (auth gate, tabs, modals); each feature exports its
        screens or a nested stack + typed \ct{ParamList}; the shell composes them and merges the
        features' deep-link config.
  \item \textbf{Expo Router}: routes are files in the \emph{app}'s \ct{app/}; keep them one-line
        re-exports so features stay router-agnostic and testable.
  \item Cross-feature jumps by \textbf{route name + params} (or a URL), never by importing a
        screen component.
\end{itemize}

\hd{Universal (iOS · Android · web)}
\begin{itemize}
  \item \textbf{react-native-web} renders \ct{View}/\ct{Text}/\ct{Pressable} as DOM;
        \textbf{Expo Router} gives Android, iOS and web one file-based navigation, with
        \textbf{static rendering} on web (SEO).
  \item Platform files: \ct{Map.native.tsx} / \ct{Map.web.tsx} (also \ct{.ios} / \ct{.android})
        — the bundler picks; the \emph{interface} lives in the shared file's types.
  \item Share logic and UI primitives; accept per-platform screens where UX differs (hover,
        keyboard, URL bar). Web bundle size and SSR-unsafe code (\ct{window}) are the tax.
\end{itemize}

\hd{Native modules package}
\ct{packages/native-*}: Expo Modules API (Swift/Kotlin) or a TurboModule spec + Codegen; a
typed JS facade, a \textbf{web / Jest fallback}, its own changelog. Native changes bump the app's
\ct{runtimeVersion} — OTA updates can't ship them.

\hd{Brownfield — RN inside an existing native app}
\begin{itemize}
  \item Official path: native project moves into \ct{ios/} + \ct{android/} of an RN project;
        host screens via \ct{RCTReactNativeFactory} (iOS) / \ct{ReactActivity} +
        \ct{ReactHost} (Android). One shared runtime; each RN screen is an
        \ct{AppRegistry} name.
  \item At scale: ship RN as a \textbf{prebuilt XCFramework / AAR} (e.g. Callstack's
        react-native-brownfield CLI) so native teams need no Node toolchain; native keeps the
        navigation; data via initial props + a native-module/event contract.
\end{itemize}

\hd{Traps}
\begin{itemize}
  \trap{Two \ct{react} copies $\to$ ``Invalid hook call''; two \ct{react-native} $\to$ native
        module / view registered twice. Dedupe, peer-depend.}
  \trap{A \ct{shared}/\ct{common} package everything depends on — a hidden monolith; split by
        purpose (\ct{ui}, \ct{data}, \ct{core}).}
  \trap{Barrel \ct{index.ts} re-exporting a heavy module eagerly — startup cost in every importer.}
  \trap{Native dep declared only in a feature package — not autolinked, crash at runtime.}
  \trap{Hoisting hides undeclared deps until a CI or pnpm install breaks them.}
\end{itemize}

\hd{Remember}
\textbf{Apps $\to$ features $\to$ shared $\to$ native, never sideways · exports = the API · lint
enforces · CODEOWNERS own · affected + cache · one react-native.}


\end{multicols}

\noindent{\footnotesize\color{sheetGrey}\textit{Related:} rn-state-and-data · rn-navigation ·
rn-native-modules · modularization-spm (the iOS equivalent) · lint-static-analysis}

\end{document}
